\documentclass[14pt,aspectratio=169]{beamer}

\usepackage{caption} 
\usepackage{siunitx}

% Add slides with section titles (done)
% TODO Add slide with spatially varying model
% TODO Add slides with clustering algorithm steps
% TODO Update simulation study slides with mathematical description
% TODO Fix plot sizes/height
% TODO Add QR code slide at end
% TODO Fix references slide at end
% TODO Go through text and line up with paper (may have better wording there)

% theme (you used auriga)
\usetheme{auriga}
\usecolortheme{auriga}

% Add section slides
\AtBeginSection[]{
  % \begin{frame}
  \begin{frame}[t, noframenumbering, plain]
  \vfill
  \centering
  \begin{beamercolorbox}[sep=8pt,center,shadow=true,rounded=true]{title}
    \usebeamerfont{title}\insertsectionhead\par%
  \end{beamercolorbox}
  \vfill
  \end{frame}
}

\titlegraphic{%
  
\includegraphics[height=1.5cm]{images/university-of-bath-logo.png}\hspace{1cm}%
  
\includegraphics[height=1.5cm]{images/UKRIlogo.png}%
}

\title{\LARGE Clustering of multivariate tail dependence using conditional methods}
\author[P. O'Toole]{%
  Paddy O'Toole, University of Bath \\[0.5em]
  \small \emph{Supervised by:} Christian Rohrbeck and Jordan Richards (University of Edinburgh)
}
\date{Statistics Seminar\\ \today}

\titlegraphic{%
  \centering
  
\includegraphics[height=1.5cm]{images/university-of-bath-logo.png}\hspace{1cm}%
  
\includegraphics[height=1.5cm]{images/UKRIlogo.png}%
}

\setbeamertemplate{title page}{%
  \begin{center}
    {\LARGE \bfseries \inserttitle \par} 
    \vspace{1.5em}
    {\large \insertauthor \par}
    \vspace{1em}
    {\small \insertdate \par}
    \vfill
    {\inserttitlegraphic\par}
  \end{center}
}

\begin{document}

% Plain, unnumbered title slide
\begin{frame}[plain,noframenumbering]
  \titlepage
\end{frame}

% % Three panel cartoon
\begin{frame}
  \centering
  \includegraphics[width=0.88\textwidth]{images/Weathering the Storm.png}
\end{frame}

% % Rene Magritte paintings
% \begin{frame}
%    \begin{minipage}{0.49\textwidth}
%      \begin{center}
%        \includegraphics[width=1\textwidth]{images/hegels_holiday.jpg}
%      \end{center}
%    \end{minipage}
%    \hfill
%    \begin{minipage}{0.49\textwidth}
%      \hspace*{0.8cm} \includegraphics[width=1\textwidth]{images/the_great_war.jpg} 
%    \end{minipage}
% \end{frame}

% Storm Darragh, December 2024, Bristol & Watchet, also flights from Bristol
% airport affected
\begin{frame}[t]

  \makebox[\textwidth]{%
    
\includegraphics[height=2.2cm,keepaspectratio]{images/bbc_report.png}%
  }
  \vspace{3mm}

  \begin{columns}[t,totalwidth=\textwidth]
    \column{0.49\textwidth}
      \centering
      \includegraphics[
        width=\linewidth,
        height=\textheight,
        keepaspectratio
      ]{images/fallen_tree.png}

    \column{0.02\textwidth}

    \column{0.49\textwidth}
      \centering
      \includegraphics[
        width=\linewidth,
        height=\textheight,
        keepaspectratio
      ]{images/waves.png}
  \end{columns}
\end{frame}

\begin{frame}{Problem} 
  \begin{itemize}
    \item Often want to estimate
      \[
        \mathbb{P}(X > x, Y > y) = \textcolor{red}{\mathbb{P}(Y > y \mid X > x)}\textcolor{blue}{\mathbb{P}(X > x)}
      \]
      for large $x, y$
    \item ``concomitant''/concurrent extreme events for random vector $\bf{X}$ often particularly devastating
    \item Goal: \textbf{identify trends by clustering sites with similar \textcolor{red}{tail dependence}}
  \end{itemize}
\end{frame}

\section{Dependence Modelling}

\begin{frame}[fragile]{Asymptotic dependence}
  % reserve space for your citation block
  \newlength{\TopBoxHeight}
  \setlength{\TopBoxHeight}{\dimexpr\textheight-3cm\relax}

  \vbox to \TopBoxHeight{
    \vfill
    \begin{center}
      \begin{itemize}
        % \item Asymptotic dependence is a measure of the strength of dependence between two random variables in the tail of their distribution.
        \item Coefficient of extremal dependence $\chi\in[0,1]$, 
              \[
                \chi = \lim_{u\to1}\mathbb{P}\bigl[F_1(X_1)>u\mid F_2(X_2)>u\bigr]
              \]
        \item (Increasingly strong) asymptotic dependence for $\chi > 0$.
        \item However, $\chi$ only gives summary; inference requires \textbf{dependence model}.
      \end{itemize}
    \end{center}
    \vfill
  }

  \vspace{1.5cm}
  {\scriptsize
    \rule{\linewidth}{0.4pt}\\
    \makebox[\linewidth]{
      \citeauthor{Coles1999}, \emph{\citetitle{Coles1999}} (1999)
    }
  }
\end{frame}

\begin{frame}{Ireland}
    \includegraphics[width = 0.9\linewidth]{plots/chi_map_ire.png}
\end{frame}

\begin{frame}[t, noframenumbering]{Ireland}
    \includegraphics[width = \linewidth]{plots/043_chi_elevation_pts_pres_man_crop.png}
\end{frame}

\section{Conditional extremes}

\begin{frame}{Marginal transformation}
  % allows for heteroskedastic regression model, symmetry means model is the same
  % for both positive and negative dependence, and also easy to model any
  % dependence strength
  \center
  \includegraphics[width=0.9\textwidth]{plots/transform_costelloe.png} 
\end{frame}

% Multivariate piecewise function, Z, Y independent
\begin{frame}[t]{Conditional extremes}
% Reduces multivariate model to simple univariate models!
\begin{itemize}
  \item \textbf{Heteroskedastic regression dependence model}: 
  \begin{equation*}
    \bigl(Y \mid X = x\bigr) = \alpha x + x^{\beta} Z, \text{ for } x > u
  \end{equation*}
  \vspace{-5mm}
  % with
  \begin{itemize}
    \item slope parameter $\alpha$ for $Y$ given large $X$,
    \item ``spread'' parameter $\beta \in (-\infty, 1]$ controls stochasticity of relationship between $Y$ and large $X$. 
  \end{itemize}
\end{itemize}
\end{frame}

\begin{frame}[t, noframenumbering]{Conditional extremes}
% Reduces multivariate model to simple univariate models!
\begin{itemize}
  \item <1-> \textbf{Heteroskedastic regression dependence model}: 
  \begin{equation*}
      \bigl(\bm{Y}_{-i} \mid Y_i = y \bigr) = \bm{\alpha}_{\mid i}\bm{y} + \bm{y}^{\bm{\beta}_{\mid i}}\bm{Z}_{\mid i}, \text{ for } y > u_{i}
  \end{equation*}
  \vspace{-5mm}
  % with
  \begin{itemize}
    \item <1-> slope parameter $\alpha_{j \mid i} \in [-1, 1]$  for $Y_j$ given large $Y_i$,
    \item <1-> ``spread'' parameter $\beta_{j \mid i} \in (-\infty, 1]$ controls stochasticity of relationship between $Y_j$ and large $Y_i$. 
  \end{itemize}
  \vspace{2mm}
  \item <2-> Key assumptions:
  \begin{itemize}
    \item <2-> Residuals $\bm{Z}_{\mid i} \sim N(\bm{\mu}_{\mid i}, \bm{\Sigma}_{\mid i})$ 
    \item <2-> $\bm{Z}_{\mid i}, Y_i$ conditionally independent for large $Y_i$
  \end{itemize}
  \vspace{2mm}
\end{itemize}
\vspace{1.3cm}
{\scriptsize
  \rule{\linewidth}{0.4pt}\\
  \makebox[\linewidth]{
    \citeauthor{Heffernan2004}, \emph{\citetitle{Heffernan2004}} (2004)
  }
}
\end{frame}

\begin{frame}[t, noframenumbering]{Conditional extremes}
% Reduces multivariate model to simple univariate models!
\begin{itemize}
  \item <1-> \textbf{Heteroskedastic regression dependence model}: 
  \begin{equation*}
    \bigl(\bm{Y}_{-i, s} \mid Y_{i, s} = y \bigr) = \bm{\alpha}_{\mid i, s}\bm{y} + \bm{y}^{\bm{\beta}_{\mid i, s}}\bm{Z}_{\mid i, s}, \text{ for } y > u_{i, s}
  \end{equation*}
  \vspace{-5mm}
  % with
  \begin{itemize}
    \item <1-> slope parameter $\alpha_{j \mid i, s} \in [-1, 1]$  for $Y_{j, s}$ given large $Y_{i, s}$,
    \item <1-> ``spread'' parameter $\beta_{j \mid i, s} \in (-\infty, 1]$ controls stochasticity of relationship between $Y_{j, s}$ and large $Y_{i, s}$. 
  \end{itemize}
  \vspace{2mm}
  \item <1-> Key assumptions:
  \begin{itemize}
    \item <1-> Residuals $\bm{Z}_{\mid i, s} \sim N(\bm{\mu}_{\mid i, s}, \bm{\Sigma}_{\mid i, s})$ 
    \item <1-> $\bm{Z}_{\mid i, s}, Y_{i, s}$ conditionally independent for large $Y_{i, s}$
  \end{itemize}
  \vspace{2mm}
\end{itemize}
\vspace{1.3cm}
{\scriptsize
  \rule{\linewidth}{0.4pt}\\
  \makebox[\linewidth]{
    \citeauthor{Heffernan2004}, \emph{\citetitle{Heffernan2004}} (2004)
  }
}
\end{frame}

\begin{frame}{Conditional extremes}
  \center{}
  \includegraphics[width=0.7\textwidth, keepaspectratio]{plots/041_costelloe_fishery.png} 
\end{frame}

\begin{frame}{Inference}
  \begin{itemize}
    Inference assumes conditional distribution follows a multivariate Normal (MVN) distribution: \\
      \begin{equation*} \label{eq:ce_norm}
        % \bigl( \bm{Y}_{-i} \mid Y_i = y_i \bigr) \sim N\left(\bm{\alpha}_{\mid i} y_i + y_i^{\bm{\beta}_i} \bm{\mu}_{\mid i}, y_i^{\bm{\beta}_i} \bm{\Sigma_{\mid i}}\right), \text{ for } Y_i > u_{Y_i}
        \bigl( \bm{Y}_{-i, s} \mid Y_{i, s} = y \bigr) \sim N\left(\bm{\alpha}_{\mid i, s} y + y^{\bm{\beta}_{\mid i, s}} \bm{\mu}_{\mid i, s}, y^{\bm{\beta}_{\mid i, s}} \bm{\Sigma_{\mid i, s}}\right), \text{ for } Y_{i, s} > u_{i, s}
      \end{equation*} \\
    $\implies$ \textbf{dependence structures at different sites can be compared using their MVN distributions}
\end{itemize}
\end{frame}

\section{Clustering}
% TODO Fill in!

\section{Simulations}

\begin{frame}{Gaussian copula}
  \begin{itemize}
    \item \num[group-separator={,}]{10000} observations of $d=2$ variables at $D = 12$ ``sites''
    \item Bivariate Gaussian copula with correlation $\rho^{s}_{\text{Gauss}} \in [0, 1]$ at site $s = 1, \ldots, D$.
    \item (asymptotic) CE parameters are $\alpha^{s} = (\rho^{s}_{\text{Gauss}})^2, \beta^{s} = 1/2$
    \item Example: 3 clusters of 4 locations defined by respective $\rho^{s}_{\text{Gauss}}$ values of $0.1, 0.5, 0.9$.
  \end{itemize}

\end{frame}

\begin{frame}{Gaussian copula}
  \center
  \includegraphics[width=0.75\textwidth, keepaspectratio]{plots/gauss_sim_ex_man.png}
\end{frame}

\begin{frame}{Gaussian copula - results}
  \center
  \includegraphics[width=0.75\textwidth, keepaspectratio]{plots/sim_01_gauss_cop_man.png}
\end{frame}

\begin{frame}{Mixture simulation}
  \begin{itemize}\setlength{\itemsep}{2pt}  % smaller space (default ~6pt)
    \item <1-> Mixture of Gaussian and t-copulas, \\
      inducing asymptotic independence and dependence, respectively
    \item <2-> \num[group-separator={,}]{10000} observations of $d=2$ variables at $D = 12$ ``sites''
    \item <2-> 2 clusters each of 6 sites
    \item <3-> Grid search performed over $\rho^s_{\text{Gauss}}, \rho^s_{t} = \rho_{t_1}, \rho^s_{t} = \rho_{t_2} \in [0, 1]$.
    \item <4-> Clustering compared to competing Vignotto algorithm using Adjusted Rand Index $\text{ARI} \in [-0.5, 1]$ 
  \end{itemize}

{\only<4>{\scriptsize
    \rule{\linewidth}{0.4pt}\\
    \begin{minipage}{\linewidth}
      \vspace{2mm}
      \citeauthor{Vignotto2021},\\
      \emph{\citetitle{Vignotto2021}} (2021)
    \end{minipage}
  }}
\end{frame}

% \begin{frame}{Mixture simulation}
%   \center
%   \includegraphics[width=1\textwidth]{plots/t_sim_ex.png}
% \end{frame}

\begin{frame}{Mixture simulation}
  \center
  \includegraphics[width=0.8\textwidth, keepaspectratio]{plots/sim_01b_ce_vs_vi_dqu_0.9_man.png}
\end{frame}

\begin{frame}{Mixture simulation}
  \center
  \includegraphics[width=0.8\textwidth, keepaspectratio]{plots/sim_01d_js_sens_3_var_dqu_0.9mult_var_man.png}
\end{frame}

\section{Ireland meteorological data}

% TODO Move colour schemes to the right of plots, can fit in more
\begin{frame}{$\alpha_{i, s}, \beta_{i, s}$ estimates}
  \vspace{-0.2cm}
  \center
  \includegraphics[width=0.6\textwidth, keepaspectratio]{plots/ire_ce_new.png} 
\end{frame}

% Maps of clustering
\begin{frame}{Ireland clustering}
  \center
  \vspace{-1cm}
  \includegraphics[width=0.9\textwidth]{plots/cluster_dqu_sep_new_crop.png} 
\end{frame}

% Sensitivity to choice of dependence quantile
\begin{frame}{Sensitivity Analysis}
  \center
  % \vspace{-0.5cm}
  \vspace{-0.5cm}
  \includegraphics[width=0.9\textwidth, keepaspectratio]{plots/cluster_dqu_crop.png} 
\end{frame}

% Dissimilarity matrix
\begin{frame}{Dissimilarity Matrix}
  \center
  % \includegraphics[width=0.75\textwidth, keepaspectratio]{plots/ire_dist_mat_heatmap_man.png} 
  % \vspace{-0.5cm}
  \vspace{-0.3cm}
  \includegraphics[width=0.65\textwidth, keepaspectratio]{plots/ire_dist_heatmap_no_labs_crop.png} 
\end{frame}

% Results for k = 2 and k = 4
% TODO Add results for k = 3 here as well for context
\begin{frame}{$k = 2$ and $k = 4$ clusters}
  \center
  \includegraphics[width=0.75\textwidth, keepaspectratio]{plots/cluster_dqu_k_2_4.png}
\end{frame}

\section{Discussion}

\begin{frame}{Discussion}
  \begin{itemize}
    \item <1-> \textbf{Conclusion}:
      \begin{itemize}
        \item Principled \& effective clustering framework for CE models
        \item Simulations \& application show clustering can group sites with similar tail dependence, which may aid interpretation
      \end{itemize} 
    \vspace{2mm}
    \item <2-> \textbf{Limitations}:
      \begin{itemize}
        \item \textbf{CE}: Gaussian assumption for $\bm{Z}_{\mid i}$
        \item \textbf{Clustering}: Uncertainty in CE fits is ignored in clustering
      \end{itemize} 
  \end{itemize} 
\end{frame}

% TODO QR Code slide

\begin{frame}\frametitle{References}
\tiny
\printbibliography[heading=none]
\end{frame}

\section{Appendix}

% TODO Add elbow & threshold stability plots for Ireland data

% \begin{frame}{Mixture simulation}
%   \center
%   \includegraphics[width=1\textwidth]{plots/t_sim_ex.png}
% \end{frame}

\begin{frame}{Mixture simulation (three variables)}
  \center
  \includegraphics[width=0.75\textwidth, keepaspectratio]{plots/sim_01d_js_sens_3_var_dqu_0.9_man.png}
\end{frame}

\begin{frame}{More realistic example}
  \center
  \includegraphics[width=0.8\textwidth, keepaspectratio]{plots/01c2_js_sens_sim_60_locs_dqu_0.9_smooth_man.png}
\end{frame}

\end{document}
